\par\noindent\SourceLeaf{53}%
and whose general term $h^{\text{tus}}$ is

\[
= \frac{1}{2^{2q+2h-2}} {}^{2q+2h-1}\mathfrak{B}^{(2q+h-1)}  .  {}^{\nu}\mathfrak{B}^{(2q+2h-1)} n^{2q+2h-1}.
\]

C) The law for the coefficient belonging to the cosine of an even multiple, for example $\cos 2q\varphi$, is found in the same way. For the cosine of such a multiple occurs only in even powers of $\cos\varphi$, starting with the one whose exponent is $2q$. The initial term of the series that expresses the coefficient of $\cos 2q\varphi$ will therefore be taken from the term of the fundamental series that contains $\cos\varphi^{2q}$, that is, from ${}^{\nu}\mathfrak{B}^{(2q)} n^{2q}\cos\varphi^{2q}$, and specifically from its $q^{\text{to}}$ term. Thus it becomes (since the expansion of $\cos\varphi^{2q}$ has as its $q^{\text{tus}}$ term

\[
= \frac{1}{2^{2q-1}} {}^{2q}\mathfrak{B}^{(2q)} \cos 2q\varphi = \frac{1}{2^{2q-1}} \cos 2q\varphi) \qquad = \frac{1}{2^{2q-1}} {}^{\nu}\mathfrak{B}^{(2q)} n^{2q}.
\]

\noindent In general, the $h^{\text{to}}$ even term of the fundamental series after this first one, and hence the quantity ${}^{\nu}\mathfrak{B}^{(2q+2h)} n^{2q+2h} \cos\varphi^{2q+2h}$, gives the coefficient belonging to $\cos 2q\varphi$ its $h^{\text{ta}}$ part. Now when $\cos\varphi^{2q+2h}$ is expanded, its $q^{\text{tus}}$ term is

\[
= \frac{1}{2^{2q+2h-1}} {}^{2q+2h}\mathfrak{B}^{(2q+h)} \cos 2q\varphi.
\]

\noindent Thus, for the coefficient belonging to $\cos 2q\varphi$, the $h^{\text{tus}}$ term after the first is

\[
= \frac{1}{2^{2q+2h-1}} {}^{2q+2h}\mathfrak{B}^{(2q+h)}  .  {}^{\nu}\mathfrak{B}^{(2q+2h)} n^{2q+2h}.
\]

D) The task we set ourselves is indeed complete. For we have duly exhibited the initial term (cf. A) of the series obtained by putting $(1+n\cos\varphi)^{\nu}$ into order after substituting a series of cosines of multiples for each power. We have also duly exhibited the initial and general terms of the appropriate coefficient for both kinds of terms that can occur in the series. One kind contains the cosine of an odd multiple, $\cos(2q-1)\varphi$ (cf. B), and the other the cosine of an even multiple, $\cos 2q\varphi$ (cf. C). But this description of the law we have found can be shortened by observing that exactly the same law holds for the coefficients of the cosines of even multiples as for those of odd multiples. This can be seen as follows. The initial term of the coefficient belonging to $\cos(2q-1)\varphi$ was

\[
\frac{1}{2^{2q-2}} {}^{\nu}\mathfrak{B}^{(2q-1)} n^{2q-1},
\]

\noindent and its general term $h^{\text{tus}}$ was

\[
\frac{1}{2^{2q+2h-2}} {}^{2q+2h-1}\mathfrak{B}^{(2q+h-1)}  .  {}^{\nu}\mathfrak{B}^{(2q+2h-1)} n^{2q+2h-1}.
\]
